\documentclass[11pt,a4paper]{amsart}

\usepackage[margin=1in]{geometry}
\usepackage{amsmath,amssymb,amsthm,mathtools}
\usepackage[expansion=false]{microtype}
\usepackage{enumitem}
\usepackage[colorlinks=true,linkcolor=black,citecolor=black,urlcolor=black]{hyperref}
\usepackage{cleveref}
\hypersetup{
  pdftitle={Cofactor independence and dilation for skew Schur functions},
  pdfauthor={John Shields},
  pdfkeywords={skew Schur function, Jacobi--Trudi determinant, cofactor independence,
    dilation, diagram reconstruction, beta numbers, skew-equivalence}
}
\pdfinfo{ /ORCID (https://orcid.org/0009-0004-3882-4571) }

\numberwithin{equation}{section}
\theoremstyle{plain}
\newtheorem{theorem}{Theorem}[section]
\newtheorem{proposition}[theorem]{Proposition}
\newtheorem{lemma}[theorem]{Lemma}
\newtheorem{corollary}[theorem]{Corollary}
\theoremstyle{definition}
\newtheorem{definition}[theorem]{Definition}
\newtheorem{example}[theorem]{Example}
\newtheorem{remark}[theorem]{Remark}
\newtheorem{question}[theorem]{Question}

\newcommand{\Q}{\mathbb Q}
\newcommand{\Z}{\mathbb Z}
\newcommand{\K}{\mathcal K}

\title{Cofactor independence and dilation\\
for skew Schur functions}
\author{John Shields}
\date{September 12, 2026}
\thanks{Independent researcher, Ginza 2-15-2, KR Ginza II 5F, Chuo-ku, Tokyo, Japan 104-0061.\\
\hspace*{\parindent}Corresponding author: \href{mailto:math@johnnyshields.com}{math@johnnyshields.com}.  ORCID: \href{https://orcid.org/0009-0004-3882-4571}{0009-0004-3882-4571}.\\
\hspace*{\parindent}This work is archived at \href{https://doi.org/10.5281/zenodo.22732784}{doi:10.5281/zenodo.22732784}.}
\subjclass[2020]{Primary 05E05; Secondary 15A15}
\keywords{skew Schur function, Jacobi--Trudi determinant, cofactor independence, dilation, diagram reconstruction, beta numbers, skew-equivalence}

\begin{document}

\begin{abstract}
The signed cofactors of a Jacobi--Trudi matrix with strictly positive indices are linearly independent over the rationals unless its skew diagram is fixed by rotation.  In the symmetric case, reflection gives every relation and an explicit independent half of the cofactors.  We prove this in every matrix order, by first reconstructing diagrams with nonnegative Jacobi--Trudi indices from their skew Schur functions.  We then determine how the first-variation kernel changes under integral dilation of the beta numbers.  On the nonsymmetric dense class the dilated kernel is the direct sum of the inherited kernel and an explicitly described arithmetic space supported off the original grid.  A finite-word period argument gives its dimension without determinant expansion.  On symmetric dense cores the kernel is determined outright.  The reconstruction theorem concerns diagram rigidity, the dilation formula linear restoration data; an explicit unequal pair shows why the latter must not be identified with full skew Schur equality.
\end{abstract}

\maketitle

\section{Introduction}

Repeated entry labels in a Jacobi--Trudi matrix are the only candidate source
of linear relations among its cofactors.  In the all-positive case they
produce none beyond rotation of the diagram.  This gives a criterion for
independence in the beta lists \(A,B\) below, with no genericity assumption on
the differences of their entries and no bound on the matrix order.

Work over \(\Q[h_1,h_2,\ldots]\), with \(h_0=1\) and \(h_j=0\) for
\(j<0\).  Let
\[
 A=(N=a_1>\cdots>a_n),\qquad
 B=(b_1>\cdots>b_n=0),\qquad a_n>b_1.
\]
These lists specify a connected skew diagram \(C\), up to translation, whose
Jacobi--Trudi matrix and signed cofactors are
\[
 M_C=(h_{a_i-b_j})_{i,j=1}^n,
 \qquad D_{ij}=(-1)^{i+j}\det M_C(\widehat i,\widehat j).
\]
The condition \(a_n>b_1\) says that every index is positive; we call the
matrix \(M_C\) and the diagram \(C\) \emph{dense}.  It is stronger than
connectedness.  Rotation through \(180\) degrees, denoted by \(C^*\), fixes
the translation class precisely when
\begin{equation}\label{eq:rotation-condition}
 a_i+b_{n+1-i}=N\qquad(1\le i\le n).
\end{equation}

\begin{theorem}[Cofactor relations]\label{thm:cofactor-relations}
Let \(C\) have the strict normalized beta lists above, with \(n\ge2\) and
\(a_n>b_1\).  If \(C\ne C^*\), its \(n^2\) signed cofactors are linearly
independent over \(\Q\).  If \(C=C^*\), then
\[
 \left\{(c_{ij}):\sum_{i,j}c_{ij}D_{ij}=0\right\}
 =\operatorname{span}_{\Q}
 \{e_{ij}-e_{n+1-j,n+1-i}:1\le i,j\le n\}.
\]
Here \(e_{ij}\) is the coordinate vector at site \((i,j)\).  In particular,
\[
 \dim_{\Q}\operatorname{span}\{D_{ij}\}
 =\begin{cases}
 n^2,&C\ne C^*,\\
 n(n+1)/2,&C=C^*.
 \end{cases}
\]
\end{theorem}

The symmetry clause concerns relations at specified matrix sites, not merely
a rank count.  It also supplies a basis: take one cofactor from each orbit of
\((i,j)\mapsto(n+1-j,n+1-i)\).  The nonsymmetric clause holds however
many entry labels repeat.  A degree count separates sites with distinct
entry labels, so only the repeated ones can carry a relation.  Both clauses
turn on extremal monomials and private labels: a square-weight ordering makes
the leading monomials of the symmetric cofactors distinct, and differentiating
at a label that occurs at a single site reduces the nonsymmetric case to a
smaller diagram.

The geometric input is a rigidity theorem.  A connected skew diagram whose
full Jacobi--Trudi matrix has no negative indices is determined, up to
translation and rotation, by its skew Schur function
(\Cref{thm:nonnegative-rigidity}).  A diagram with the same skew Schur
function is not required to satisfy that index condition.  The coefficient of
the largest label identifies a smaller diagram; the remaining polynomial
places the two exterior beta gaps.

The application of \Cref{thm:cofactor-relations} is to common beta entries.
Removing the intersection of the two beta lists leaves a \emph{zero-free
core}, in which no entry lies in both lists, so that no matrix entry is
\(h_0\).
Restoring a common entry at \(z\) contributes, to first order,
\begin{equation}\label{eq:unit-intro}
 U_z^C=\sum_{a_i>z>b_j}D_{ij}h_{a_i-z}h_{z-b_j}.
\end{equation}
Write \(\K(C)\), the \emph{first-variation kernel}, for the rational relation
space of these polynomials, indexed by the available positions between \(0\)
and \(N\).
For a positive integer \(t\), dilation replaces \((A,B)\) by \((tA,tB)\).
\Cref{thm:dilation} proves that \(\K(tC)\) splits into the old kernel on the
positions divisible by \(t\) and an arithmetic kernel on the remaining
positions, whenever the signed cofactors of \(C\) are independent.  By
\Cref{thm:cofactor-relations} that holds for every nonsymmetric dense core, so
there the formula carries no hypothesis beyond density and nonsymmetry.  The
arithmetic part is given explicitly by odd functions on a finite residue
group.  Thus the new relations introduced by dilation can be determined
without knowing a formula for the entire base kernel.

Skew Schur equality is a broader problem.  McNamara and van Willigenburg
conjecture a classification by composition and reversal of factors
\cite{McNamara2009Classification}; for ribbons, the corresponding
classification is due to Billera, Thomas and van Willigenburg
\cite{Billera2006DecomposableCompositions}.
Neither a linear kernel nor a finite list of equal diagrams supplies the
factor presentations required by the necessity direction of that
classification.  The gap between restoration data and full equality remains
even when the core is dense and all its cofactors are independent:
\Cref{ex:unequal-restorations} gives two restorations whose lowest two ordinary
homogeneous components agree but whose full functions differ.

The arguments below use no finite census or rank certificate as a premise.
The external algebraic input is irreducibility of a connected skew Schur
function, in the corrected form proved by Barekat, Reiner and van Willigenburg
\cite[Thm.~1]{Barekat2009CoincidencesCorrigendum}.  For the standard determinant,
rotation, and row-overlap facts see
\cite{Macdonald1995SymmetricFunctions,Reiner2007Coincidences}.
Section~\ref{sec:boundary} proves the geometric rigidity theorem, and
Section~\ref{sec:cofactors} proves \Cref{thm:cofactor-relations}.
Section~\ref{sec:restoration} derives the first-variation polynomials.
Section~\ref{sec:dilation} proves the dilation formula, and
Section~\ref{sec:full-equality} explains its consequences for full equality.

\section{Boundary coefficients and dense rigidity}\label{sec:boundary}

A connected diagram whose Jacobi--Trudi indices are all nonnegative is
determined, up to translation and rotation, by its skew Schur function.
That rigidity is the geometric input
to both \Cref{thm:cofactor-relations} and \Cref{cor:independent-restorations}.

\subsection{Beta lists and polynomial readings}

Use one-based row indices.  For a skew diagram \(D=\lambda/\mu\) with
\(n\) nonempty rows, put
\[
 x_i=\lambda_i-i,\qquad y_i=\mu_i-i.
\]
A common translation of both lists changes no matrix entry.  We normalize by
\(y_n=0\), and put \(N=x_1\).  The row lengths and adjacent overlaps minus
one are
\begin{equation}\label{eq:profile}
 r_i=x_i-y_i>0,\qquad
 s_i=x_{i+1}-y_i,\qquad
 0\le s_i<\min(r_i,r_{i+1}).
\end{equation}
The last inequalities characterize connectedness.  Conversely, any ordered
profile satisfying \eqref{eq:profile} determines a connected diagram up to
translation.  Starting with \(y_n=0\) and \(x_n=r_n\), recover its lists by
\[
 y_i=x_{i+1}-s_i,\qquad x_i=y_i+r_i
 \qquad(i=n-1,\ldots,1).
\]
Rotation reverses both profile words.  On normalized beta lists it sends
\((X,Y)\) to
\[
 \bigl((N-y_n,\ldots,N-y_1),
       (N-x_n,\ldots,N-x_1)\bigr).
\]
Reversing rows and columns and transposing the determinant proves
\(s_{D^*}=s_D\).

The weighted degree of \(h_j\) is \(j\); its \emph{ordinary degree} is one
for \(j>0\).  These two gradings will have different uses.  For a monomial
\(m=\prod h_{d_i}\), its square weight is \(\sum d_i^2\).
Zeros may be included in this sum although \(h_0=1\).

\begin{lemma}\label{lem:polynomial-readings}
Let \(P=s_D\) for a nonempty connected skew diagram with full beta lists
\(X,Y\).  Its maximum ordinary degree is the number \(n\) of rows.  Its
unique minimum-square monomial has coefficient one and reads the row-length
multiset.  The largest generator index occurring in \(P\) is \(N=x_1\).
Moreover,
\begin{equation}\label{eq:corner-coefficient}
 Q=(-1)^{n-1}[h_N]P
   =\det(h_{x_{i+1}-y_j})_{i,j=1}^{n-1}\ne0.
\end{equation}
The unique minimum-square monomial of \(Q\), padded with zeros to length
\(n-1\), reads the multiset of the \(s_i\).
For \(n=1\), the determinant on the right is one.
\end{lemma}

\begin{proof}
For a permutation \(\pi\),
\[
 \sum_i(x_i-y_{\pi(i)})^2
 =\sum_i x_i^2+\sum_i y_i^2-2\sum_i x_i y_{\pi(i)}.
\]
Strict rearrangement makes the identity permutation the unique minimizer.
Its diagonal indices are positive.  No other determinant term can give the
same monomial, since it would have the same square weight.  The diagonal
therefore survives with coefficient one and ordinary degree \(n\); no term
has larger ordinary degree.

The label \(N\) occurs at the unique site \((1,n)\).  Its complementary
board has strict lists and nonnegative diagonal \((s_1,\ldots,s_{n-1})\).
The same rearrangement argument proves its nonvanishing and reads that
diagonal, allowing zeros.  Hence the corner label really occurs in \(P\),
and no matrix label is larger.
\end{proof}

The readings in \Cref{lem:polynomial-readings} are standard consequences of
Jacobi--Trudi; see \cite{Reiner2007Coincidences}.  They assume neither density
nor upper Hessenberg form, so they are available on an unknown equal diagram
before its shape has been recovered.

We also use the following established fact.  A nonempty connected skew Schur
function is irreducible in \(\Z[h_1,h_2,\ldots]\)
\cite[Thm.~1]{Barekat2009CoincidencesCorrigendum}, hence in
\(\Q[h_1,h_2,\ldots]\) by primitivity.  A disconnected skew diagram has
function equal to the product of those of its nonempty connected components.
Consequently an irreducible skew Schur function can have only connected
realizations after empty rows are removed.  This is algebraic irreducibility,
not irreducibility under composition.

\subsection{Placing the exterior gaps}

For a diagram \(T\) with strict lists \(U,V\), write
\(T^{\mathrm{top}}(d)\) for the diagram obtained by increasing its largest
row beta by \(d>0\), and \(T^{\mathrm{bottom}}(d)\) for the diagram obtained
by decreasing its smallest column beta by \(d\).  These operations extend
the first and last rows, respectively.  The \emph{maximal corner cofactor} of
a diagram is the one at the site carrying the largest label, namely
\((1,n)\); its lists are \((x_2,\ldots,x_n)\) and \((y_1,\ldots,y_{n-1})\),
and its polynomial is \eqref{eq:corner-coefficient}.

\begin{lemma}\label{lem:exterior-gaps}
Let a connected diagram have \(n\ge3\) rows and a nonnegative-index
Jacobi--Trudi matrix.  Write the strict lists of its maximal corner cofactor
as
\[
 U=(u_1>\cdots>u_m),\quad V=(v_1>\cdots>v_m),
 \qquad m=n-1.
\]
Its full lists have the form
\begin{equation}\label{eq:border}
 X=(u_1+\alpha,u_1,\ldots,u_m),\qquad
 Y=(v_1,\ldots,v_m,v_m-\beta),\qquad \alpha,\beta>0.
\end{equation}
Put \(N=u_1-v_m+\alpha+\beta\).
The largest generator index in \(s_D|_{h_N=0}\) is \(M=N-\min(\alpha,\beta)\).
If \(\alpha<\beta\), its coefficient is
\((-1)^n s_{T^{\mathrm{top}}(\alpha)}\); if \(\beta<\alpha\), it is
\((-1)^n s_{T^{\mathrm{bottom}}(\beta)}\).
If \(\alpha=\beta\), the coefficient of \(h_M^2\) is nonzero.
\end{lemma}

\begin{proof}
Apart from the corner, the two largest possible indices occur at
\((2,n)\) and \((1,n-1)\), with values \(N-\alpha\) and \(N-\beta\).
If the values differ, the larger is private.  Deleting its row and column
gives exactly the stated extension of \(T\), with sign \((-1)^n\).
That determinant is nonzero and contains no \(h_N\), so the specialization
at \(h_N=0\) does not remove it.

If the values agree, both sites can be used simultaneously.  They are the
only sites carrying \(M\): here \(M=x_2-y_n\), so a site \((i,j)\) carries it
exactly when \(x_i-x_2=y_j-y_n\), which fails for \(i>2\) because the left
side is negative and the right side is not, and for \(i=1\) forces
\(y_j=y_n+\alpha=y_{n-1}\).  Their complementary board has row lists
\((x_3,\ldots,x_n)\) and column lists \((y_1,\ldots,y_{n-2})\), with
nonnegative diagonal.  Strict rearrangement proves its determinant nonzero.
This is, up to sign, the coefficient of \(h_M^2\).  No other remaining site
has a larger index.
\end{proof}

\begin{theorem}[Nonnegative-index rigidity]\label{thm:nonnegative-rigidity}
Let \(D\) be connected with \(n\) nonempty rows and full beta lists \(X,Y\).
If \(x_n\ge y_1\), then every connected diagram \(D'\) with \(s_{D'}=s_D\)
is a translate of \(D\) or \(D^*\).
\end{theorem}

\begin{proof}
Induct on \(n\).  One row is immediate.  In two rows, the row multiset and
the single adjacent overlap determine the diagram up to rotation, so
\Cref{lem:polynomial-readings} proves the assertion even without the
nonnegative-index assumption.

Suppose \(n\ge3\).  The maximal corner cofactor \(T\) of \(D\) has
\(n-1\) positive-length rows, is connected, and has no negative indices.
Indeed its least index is still \(x_n-y_1\ge0\); its diagonal is strictly
positive because the full lists are strict and \(n\ge3\).
Its polynomial is irreducible.

By \Cref{lem:polynomial-readings}, \(D'\) has the same row count and the same
signed corner coefficient \eqref{eq:corner-coefficient}.  Its padded diagonal reading
shows that the corresponding cofactor of \(D'\) also has \(n-1\)
positive-length rows.  Irreducibility makes it connected.  Induction therefore
identifies it with \(T\) or \(T^*\).  Rotate and translate \(D'\) to align
these cofactor lists.

Both parents now have form \eqref{eq:border}.  In particular the least index
of the mate is \(u_m-v_1\), the same as that of the source.  Thus it has
entered the nonnegative-index class; this was not assumed.
The common largest label determines \(\alpha+\beta\), and
\Cref{lem:exterior-gaps} determines \(\min(\alpha,\beta)\).  Only interchange
of the two margins remains.

If the margins agree there is nothing to decide.  Otherwise put
\(d=\min(\alpha,\beta)\).  An interchange preserving the full polynomial
would give
\[
 s_{T^{\mathrm{top}}(d)}=s_{T^{\mathrm{bottom}}(d)}.
\]
Both extensions have \(n-1\) rows and remain in the nonnegative-index class.
Induction says they agree up to translation or rotation.  They cannot be
translates: their largest column beta agrees, forcing translation zero,
whereas their largest row betas differ.  If they are rotations, the reflection
constant is \(u_1+v_m\).  Undoing the two endpoint changes in their beta lists
gives
\(u_i+v_{m+1-i}=u_1+v_m\) for \(1\le i\le m\).
Thus \(T=T^*\), and the two parents with interchanged margins are themselves
rotations.  This completes the induction.
\end{proof}

\section{Linear relations among cofactors}\label{sec:cofactors}

Every cofactor of a dense matrix is, up to its specified sign, the skew Schur
function of a connected dense diagram.  In particular it is nonzero.
Its weighted degree is
\begin{equation}\label{eq:cofactor-degree}
 \deg D_{ij}=|C|-(a_i-b_j),\qquad
 |C|=\sum_i a_i-\sum_j b_j.
\end{equation}
A rational relation therefore splits into groups of sites carrying the same
entry label.  Repeated labels are the only groups requiring attention.

\subsection{The symmetric case}

\begin{proof}[Proof of the symmetric clause of \Cref{thm:cofactor-relations}]
Assume \eqref{eq:rotation-condition}.  Reverse the columns of \(M_C\).
The resulting symmetric matrix is \(S_{ij}=h_{a_i+a_j-N}\).
Put
\[
 d_i=2a_i-N,\qquad e_i=a_i+a_{i+1}-N.
\]
Density and symmetry give
\begin{equation}\label{eq:interlaced-indices}
 d_1>e_1>d_2>e_2>\cdots>e_{n-1}>d_n>0.
\end{equation}

For \(i\le j\), delete row \(i\) and column \(j\) from \(S\).
Strict rearrangement now makes the \emph{maximum}-square matching unique:
the two remaining decreasing lists of \(a\)'s are paired in the same order.
Its monomial is
\begin{equation}\label{eq:symmetric-leader}
 h_{d_1}\cdots h_{d_{i-1}}
 h_{e_i}\cdots h_{e_{j-1}}
 h_{d_{j+1}}\cdots h_{d_n}.
\end{equation}
All its indices are positive, and its coefficient is a unit.
The monomials in \eqref{eq:symmetric-leader} are distinct for distinct pairs
\(i\le j\): the strictly interlacing list
\eqref{eq:interlaced-indices} records exactly which consecutive string of
\(e\)'s replaces the omitted \(d\)'s.

Order monomials by square weight and then by any fixed total tie-break.
These cofactors have distinct leading monomials, hence are independent.
Symmetry of \(S\) identifies the cofactor at \((i,j)\) with that at
\((j,i)\).  These are all the relations: the \(n(n+1)/2\) cofactors with
\(i\le j\) are independent, so the relation space has dimension
\(n^2-n(n+1)/2=n(n-1)/2\), which the differences already exhibited attain.
Undoing the column reversal changes
the site involution to
\((i,j)\mapsto(n+1-j,n+1-i)\).
\end{proof}

\subsection{The nonsymmetric case}

\begin{lemma}\label{lem:boundary-pair}
Let \(C\) be nonsymmetric and dense, with normalized lists \(A,B\).
Suppose \(j<n\), \(i>1\), and \(a_i+b_j=N\).
Then \(D_{1j}\) and \(D_{in}\) are not proportional over \(\Q\).
\end{lemma}

\begin{proof}
Their unsigned cofactor diagrams have lists
\[
 E=(A\setminus\{N\},B\setminus\{b_j\}),\qquad
 F=(A\setminus\{a_i\},B\setminus\{0\}).
\]
If their polynomials are proportional, their coefficient-one
minimum-square monomials make the unsigned scalar one.
Both diagrams are dense and connected, so
\Cref{thm:nonnegative-rigidity} makes them translates or rotations.

For a rotation from \(E\) to \(F\), the reflection constant is \(N\),
since the largest row beta of \(F\) is \(N\) and the smallest column beta
of \(E\) is zero.  Restoring the omitted pair
\(a_i=N-b_j\) then gives \(A=N-B\) as sets, contrary to nonsymmetry.

For a translation from \(E\) to \(F\), comparison of extrema gives the
positive displacement \(u=N-a_2=b_{n-1}\).
The minimum row beta forces \(i=n\), since otherwise both row lists have
minimum \(a_n\).  The maximum column beta similarly forces \(j=1\).
The remaining equalities say
\[
 (a_2,\ldots,a_n)+u=(a_1,\ldots,a_{n-1}),\qquad
 (b_2,\ldots,b_n)+u=(b_1,\ldots,b_{n-1}).
\]
All row and column gaps are \(u\), which again implies
\eqref{eq:rotation-condition}.
\end{proof}

\begin{proof}[Proof of the nonsymmetric clause of \Cref{thm:cofactor-relations}]
Induct on \(n\).  For \(n=2\), the four matrix labels are distinct unless
\(a_2=N-b_1\), which is exactly symmetry.  The four cofactors are these
four generators up to sign.

Let \(n\ge3\), and suppose
\begin{equation}\label{eq:cofactor-relation}
 \sum_{i,j}c_{ij}D_{ij}=0,\qquad c_{ij}\in\Q.
\end{equation}
The label \(h_N\) is private at \((1,n)\).  Let
\[
 R=(A\setminus\{N\},B\setminus\{0\}).
\]
For \(i>1\) and \(j<n\), differentiation gives
\begin{equation}\label{eq:private-derivative}
 \frac{\partial D_{ij}}{\partial h_N}
   =(-1)^{n-1}D^R_{i-1,j};
\end{equation}
on the first row or last column the derivative is zero.

If \(R\ne R^*\), induction applied to \eqref{eq:private-derivative}
kills all \(c_{ij}\) with \(i>1,j<n\).  Only the first row and last
column remain.  By \eqref{eq:cofactor-degree}, each weighted-degree group
contains at most two sites, one on each of those boundaries.
A singleton cannot give a relation because its cofactor is nonzero;
a pair cannot give one by \Cref{lem:boundary-pair}.  Hence every coefficient
vanishes.

Suppose instead that \(R=R^*\).  Put \(u=N-a_2\) and \(v=b_{n-1}\).
Nonsymmetry of the parent gives \(u\ne v\).  Rotate the parent, if necessary,
to have \(u<v\).  The label \(h_{a_2}\) is then private at \((2,n)\):
every first-row label outside the corner is at most \(N-v<a_2\), and all
other candidates are smaller by strictness.

The complementary diagram \(Q=(A\setminus\{a_2\},B\setminus\{0\})\) is
nonsymmetric.  Otherwise the symmetries of \(R\) and \(Q\) would have
the same center, obtained from their common minimum row beta and maximum
column beta, but different centers obtained from their maximum row betas
\(a_2,N\) and common minimum column beta \(v\).

Differentiate \eqref{eq:cofactor-relation} at \(h_{a_2}\).
Independence for \(Q\), supplied by induction, kills every coefficient
except those on row two or column \(n\).  Now differentiate at \(h_N\)
again.  The resulting relation on the symmetric child \(R\) is supported
in its first row.  Its relation space has already been proved to be
antisymmetric under site reflection.  No nonzero vector in that space can
be supported in one row: a reflected partner lies in another row unless
the site is fixed, when its coefficient is zero.  This kills row two away
from column \(n\).

The remaining last-column cofactors have distinct weighted degrees, so their
coefficients vanish separately.  This proves the induction.
\end{proof}

Only the scalar polynomial of each smaller diagram was used to establish
\Cref{lem:boundary-pair}.  No step transports an indexed cofactor matrix
between two diagrams with equal functions.

\section{Common beta entries and first variations}\label{sec:restoration}

\subsection{The exact path expansion}

For a diagram with full beta lists \(X,Y\), set
\[
 Z=X\cap Y,\qquad A=X\setminus Z,\qquad B=Y\setminus Z,
 \qquad k=|A|=|B|,
\]
retaining decreasing order.  The pair \(C=(A,B)\) is its zero-free core;
\emph{zero-free} means \(A\cap B=\varnothing\), not that negative indices
are absent.  A dense core is zero-free, since \(a_n>b_1\) forces
\(\min A>\max B\).  The determinant still uses \(h_j=0\) for \(j<0\).

Canceling common entries preserves the difference of the two endpoint
counting functions.  For normalized connected lists this difference is at
least one at every real point strictly between \(0\) and \(N\).
Consequently a connected diagram has a connected residual core.
Conversely, for a connected zero-free core and
\[
 \Omega_C=\{1,\ldots,N-1\}\setminus(A\cup B),
\]
any subset \(Z\subseteq\Omega_C\) can be restored in both lists.
Sorting \(A\cup Z\) and \(B\cup Z\) gives a connected diagram, denoted
by \(C[Z]\).  Its area is \(|C|\), and its row count is \(k+|Z|\).

For strict lists \(U,V\), define the positive-entry matrix
\[
 H(U,V)_{uv}=\begin{cases}h_{u-v},&u>v,\\0,&u\le v.\end{cases}
\]
In this auxiliary matrix the value at equality is zero.  The original
Jacobi--Trudi matrix, in contrast, retains its genuine units.

\begin{proposition}\label{prop:path-expansion}
Let \(S=\lambda/\mu\) have \(n\) nonempty rows and the lists above.
Put
\[
 L=H(A,B),\quad P=H(A,Z),\quad Q=H(Z,B),\quad T=H(Z,Z).
\]
Let \(\epsilon\) be the product of the permutation signs for regrouping
\(X\) as \((A,Z)\) and \(Y\) as \((B,Z)\).  The variable \(t\) marks ordinary
degree, not the dilation factor of Section~\ref{sec:dilation}.
Then
\begin{equation}\label{eq:schur-complement}
 s_S(th_1,th_2,\ldots)
 =\epsilon t^k\det(L-tPQ+t^2PTQ-t^3PT^2Q+\cdots).
\end{equation}
The sum terminates.  Equivalently, the entry indexed by \(a\in A,b\in B\)
of the reduced matrix is
\begin{equation}\label{eq:path-entry}
 \sum_{\ell\ge0}(-t)^\ell
 \sum_{\substack{a>z_1>\cdots>z_\ell>b\\ z_i\in Z}}
 h_{a-z_1}h_{z_1-z_2}\cdots h_{z_\ell-b},
\end{equation}
where the term for \(\ell=0\) is \(H(a,b)\).
\end{proposition}

\begin{proof}
After regrouping, the scaled full matrix is
\[
 \begin{pmatrix}tL&tP\\tQ&I+tT\end{pmatrix}.
\]
The matrix \(T\) is strictly upper triangular.  Hence
\(\det(I+tT)=1\) and
\((I+tT)^{-1}=I-tT+t^2T^2-\cdots\).
The block determinant identity gives \eqref{eq:schur-complement}, without
requiring \(L\) to be invertible.  Matrix multiplication along the ordered
common beta entries gives \eqref{eq:path-entry}.
\end{proof}

Write \([f]_d\) for the ordinary degree-\(d\) component of \(f\).
The sorted diagonal of the residual matrix is strictly positive, so its
determinant is nonzero by strict rearrangement.  Thus the lowest ordinary
degree in \(s_S\) is exactly \(k\).  By \Cref{prop:path-expansion},
\begin{equation}\label{eq:first-components}
 [s_S]_k=\epsilon s_C,\qquad
 [s_S]_{k+1}=-\epsilon\sum_{z\in Z}U_z^C,
\end{equation}
where \(U_z^C\) is given by \eqref{eq:unit-intro}.  Every \(U_z^C\)
is homogeneous of ordinary degree \(k+1\) and weighted degree \(|C|\).
These polynomials describe insertion of one common beta entry; they are not
individual matrix entries equal to one.

\begin{definition}
For a connected zero-free core, its first-variation kernel is
\[
 \K(C)=\left\{\lambda\in\Q^{\Omega_C}:
          \sum_{z\in\Omega_C}\lambda_zU_z^C=0\right\}.
\]
\end{definition}

\subsection{Endpoint values}

Extend the formula for \(U_z^C\) to all integer positions.
Writing \(U_z^C=u^{\top}Dv\) with \(D=(D_{ij})\), the \emph{return vectors}
are \(u_i=h_{a_i-z}\) and \(v_j=h_{z-b_j}\), each entry taken as zero unless
its gap is positive.

\begin{lemma}\label{lem:endpoint-cancellation}
For every zero-free core,
\[
 U_z^C=0\qquad(z\in A\cup B).
\]
Consequently values at beta endpoints can be assigned arbitrarily in a
coefficient vector without changing its unit relation.
\end{lemma}

\begin{proof}
If \(z=a_\ell\), the right return vector is row \(\ell\) of \(M_C\).
The adjugate identity gives
\[
 \sum_jD_{ij}(M_C)_{\ell j}=\delta_{i\ell}\det M_C.
\]
The \(\ell\)-th entry of the other return vector is zero, so contraction
vanishes.  If \(z=b_\ell\), use the corresponding column identity.
Zero-freeness ensures that these matrix rows and columns have exactly the
strict-return convention required here.
\end{proof}

On a symmetric core the kernel is determined outright.

\begin{proposition}\label{prop:symmetric-unit-kernel}
Let \(C\) be a normalized symmetric dense core of order at least two.  Then
\(U_z^C=U_{N-z}^C\) for every \(z\in\Omega_C\), and
\[
 \K(C)=\{\lambda\in\Q^{\Omega_C}:\lambda_z+\lambda_{N-z}=0\}
       =\operatorname{span}_{\Q}\{e_z-e_{N-z}:z\in\Omega_C,\ 2z<N\},
\]
where \(e_z\) is the coordinate vector at \(z\).  In particular
\(\dim_{\Q}\K(C)=\lfloor|\Omega_C|/2\rfloor\).
\end{proposition}

\begin{proof}
Symmetry gives \(a_i=N-b_{n+1-i}\) and \(b_j=N-a_{n+1-j}\), so the substitution
\(i'=n+1-j\), \(j'=n+1-i\) turns \(a_i>N-z>b_j\) into \(a_{i'}>z>b_{j'}\) and the
return factors into \(h_{z-b_{j'}}h_{a_{i'}-z}\).  It is an involution, hence a
bijection of the two summation ranges.  Its coefficient is
\(D_{n+1-j',n+1-i'}\), which \Cref{thm:cofactor-relations} identifies with
\(D_{i'j'}\).  Therefore \(U_{N-z}^C=U_z^C\).  Reflection preserves \(A\cup B\),
so it preserves \(\Omega_C\), and each \(e_z-e_{N-z}\) lies in \(\K(C)\).

For the reverse inclusion it suffices that \(\{U_z^C:z\in\Omega_C,\ 2z\le N\}\)
be independent, one representative per reflection orbit.  Induct on the order.
For a one-row core \(U_z^C=h_{N-z}h_z\), and these monomials are distinct for
\(2z\le N\).

Let \(n\ge2\) and put \(u=b_{n-1}=N-a_2\), so \(0<u<N/2\).  The lists
\(A'_p=a_{p+1}-u\) and \(B'_q=b_q-u\) define a symmetric dense core \(C'\) of
order \(n-1\) and span \(N'=N-2u\).  In a relation among the representatives,
every return index is below \(N\), so differentiating at \(h_N\) acts only on the
cofactors; by \eqref{eq:private-derivative} it sends \(D_{ij}\) to
\((-1)^{n-1}D^{C'}_{i-1,j}\) for \(i>1\), \(j<n\), and to zero on the first row or
last column.  Hence \(\partial U_z^C/\partial h_N\) vanishes for \(z<u\) and
equals \((-1)^{n-1}U_{z-u}^{C'}\) for \(z>u\); the position \(u\) is a beta
endpoint and is unavailable.  Since \(z\mapsto z-u\) carries the representatives
above \(u\) bijectively onto those of \(C'\), the midpoint included, induction
kills every coefficient with \(z>u\).

Only \(z<u\) remains.  Every column beta other than \(b_n=0\) is at least \(u\),
so \(U_z^C=h_z\sum_iD_{in}h_{a_i-z}\).  The generators of \(D_{in}\) have index at
most \(N-u\), and \(a_i-z\le a_2-z<N-u\) for \(i>1\), while the term \(i=1\)
contributes \(h_{N-z}\) with \(N-z>N-u\).  Thus
\([h_{N-z}]U_z^C=h_zD_{1n}\ne0\).  Taking the smallest surviving \(z\), every
other remaining unit has all generator indices below \(N-z\), so the coefficient
of \(h_{N-z}\) in the relation is \(c_zh_zD_{1n}\), forcing \(c_z=0\).

Every relation therefore satisfies \(\lambda_z+\lambda_{N-z}=0\), which at an
available midpoint reads \(2\lambda_{N/2}=0\).  Counting one free coefficient per
two-element orbit gives the dimension.
\end{proof}

For an extended vector \(\lambda=(\lambda_z)_{1\le z<N}\), put
\[
 Q_\lambda(b,a)=\sum_{b<z<a}\lambda_z h_{a-z}h_{z-b}.
\]
For \(a>b\), this quadratic vanishes precisely when
\(\lambda_z=-\lambda_{a+b-z}\) for \(b<z<a\).
The midpoint, if present, must have coefficient zero.  Indeed two positive
gaps have the same monomial exactly when they agree as an unordered pair.

\section{Integral dilation}\label{sec:dilation}

\subsection{From edge reflections to a common period}

For a connected normalized zero-free core of order at least two, define
\begin{equation}\label{eq:endpoint-modulus}
 g(C)=\gcd\bigl(\{N-a_i\}_i\cup\{b_j\}_j\bigr)>0.
\end{equation}
Thus all row betas are \(N\) modulo \(g(C)\), and all column betas are zero.
If \(\phi:\Z/g(C)\Z\to\Q\) satisfies \(\phi(r)+\phi(N-r)=0\), then its
restriction \((\phi(z\bmod g(C)))_{z\in\Omega_C}\) lies in
\(\K(C)\).  To see this, fill in the endpoint values by the same formula,
use \Cref{lem:endpoint-cancellation}, and pair each return position with
\(a_i+b_j-z\).  This proves an arithmetic subspace of the kernel, not yet
its completeness.

A word \(w_1,\ldots,w_L\) has period \(p\) when
\(w_i=w_{i+p}\) for \(1\le i\le L-p\).
We need only the following elementary finite-word fact: periods \(p,q>0\)
with \(p+q\le L\) imply period \(\gcd(p,q)\).
We include the argument for completeness.  Assume \(p<q\).
For \(i>p\), comparison through
\(i-p\) gives period \(q-p\).  For \(i\le p\), compare through
\(i+q\), which lies in the word because \(p+q\le L\).
This proves period \(q-p\) throughout its permitted range.  Euclid's
algorithm completes the argument.

\begin{lemma}\label{lem:endpoint-normalization}
Let \(C=(A,B)\) be a connected zero-free core of order \(n\ge2\), with
\(a_1=N\) and \(b_n=0\).  A vector on \(\Omega_C\) admits endpoint values
such that
\begin{equation}\label{eq:all-edges}
 Q_\lambda(b_j,a_i)=0\qquad\text{whenever }a_i>b_j
\end{equation}
if and only if it is the restriction of a residue function modulo \(g(C)\)
satisfying \(\phi(r)+\phi(N-r)=0\).
\end{lemma}

\begin{proof}
The reverse direction was proved after \eqref{eq:endpoint-modulus}.
For the forward direction, assume \eqref{eq:all-edges} and regard the
endpoint extension as a word of length \(N-1\).  The interval \((0,N)\) gives its odd reflection.  Comparing this
reflection with that on \((0,a_i)\) gives period \(N-a_i\), and comparing
with that on \((b_j,N)\) gives period \(b_j\).  All these intervals are
positive edges: connectedness puts every row beta above zero and every column
beta below \(N\).

Start with the two positive exterior periods \(u=N-a_2\) and \(v=b_{n-1}\).
Zero-freeness and connectedness give \(a_2>b_1\) and \(a_n>b_{n-1}\), so
\(u+v\le N-1\).  The finite-word argument gives period \(d=\gcd(u,v)\).
Incorporate the remaining column periods using \(d+b_j\le u+b_1\le N-1\),
and then the remaining row periods using
\(d+(N-a_i)\le v+N-a_n\le N-1\).
The running gcd only decreases and still divides \(u,v\), so these bounds
continue to hold.  The final period is \(g(C)\).  The reflection on
\((0,N)\) gives the required oddness of the residue function.
\end{proof}

The lemma requires connectedness, not density.  It classifies vectors that
can be made to vanish on \emph{each} positive-edge quadratic.  A vanishing
sum of cofactor-weighted quadratics need not have that stronger property.
The proof of \Cref{thm:dilation} supplies it, by independence over a disjoint
polynomial alphabet.

\subsection{The direct-sum formula}

Let \(tC=(tA,tB)\), and write
\(\sigma_t(h_j)=h_{tj}\).  Dilation preserves connectedness and
zero-freeness, has span \(tN\), and has endpoint modulus \(tg(C)\).
The map \(\sigma_t\) is injective.
For \(\lambda\in\Q^{\Omega_C}\), define
\[
 (\iota_t\lambda)_{tz}=\lambda_z\quad(z\in\Omega_C),\qquad
 (\iota_t\lambda)_w=0\quad(t\nmid w).
\]
These prescriptions cover \(\Omega_{tC}\).
Define \(\mathcal A_{\mathrm{off}}(tC)\) to be the restrictions to
\(\Omega_{tC}\) of functions \(\phi:\Z/tg(C)\Z\to\Q\) satisfying
\begin{equation}\label{eq:offgrid-space}
 \phi(r)+\phi(tN-r)=0,\qquad
 \phi(r)=0\quad\text{if }t\mid r.
\end{equation}

\begin{theorem}[Dilation of the first-variation kernel]\label{thm:dilation}
Let \(C\) be a connected zero-free core of order at least two, with
linearly independent signed Jacobi--Trudi cofactors.  For every integer
\(t\ge1\),
\begin{equation}\label{eq:dilation-kernel}
 \K(tC)=\iota_t\K(C)\oplus\mathcal A_{\mathrm{off}}(tC).
\end{equation}
In particular, by \Cref{thm:cofactor-relations}, the cofactor hypothesis is
automatic for every nonsymmetric dense core.  Without that hypothesis, the
right-hand side of \eqref{eq:dilation-kernel} is still a direct sum contained
in \(\K(tC)\).
\end{theorem}

\begin{proof}
Let the core order be \(n\).  At a position divisible by \(t\),
\[
 U_{tz}^{tC}=\sigma_t(U_z^C).
\]
All \(n+1\) generator factors in each monomial have indices divisible by
\(t\).  At a position not divisible by \(t\), the \(n-1\) cofactor
factors have divisible indices and the two return factors have nondivisible
indices.  Grading by the number of nondivisible generator factors separates
any relation into a grid relation and an off-grid relation.  Injectivity of
\(\sigma_t\) identifies the first with \(\iota_t\K(C)\).

For an off-grid relation, extend its vector by zero at all positions divisible
by \(t\), including every beta endpoint.  Its full equation is
\begin{equation}\label{eq:offgrid-relation}
 \sum_{i,j}\sigma_t(D_{ij})
       Q_\lambda(tb_j,ta_i)=0,
\end{equation}
where nonpositive intervals contribute zero.
Each quadratic uses only generators with indices not divisible by \(t\).
The cofactors use only generators with indices divisible by \(t\), and
remain linearly independent over \(\Q\) by injectivity of \(\sigma_t\).
Expanding \eqref{eq:offgrid-relation} in monomials of the nondivisible
alphabet shows that every quadratic coefficient vanishes separately.

Apply \Cref{lem:endpoint-normalization} to \(tC\).
The full extended word is periodic modulo \(tg(C)\) and odd about \(tN\).
Since \(g(C)\) divides \(N-a_2>0\), it is smaller than \(N\), so every
residue of this modulus occurs in the interval
\(1,\ldots,tN-1\).  The prescribed zero values on the grid therefore give
exactly the second condition in \eqref{eq:offgrid-space}.
Conversely every such function kills each edge quadratic and hence the full
unit sum.  The two summands have disjoint supports, proving the direct sum.
\end{proof}

The hypothesis concerns cofactor independence, whereas the first summand in
\eqref{eq:dilation-kernel} retains the \emph{entire} base unit kernel.
\Cref{thm:dilation} therefore does not infer unit independence from cofactor
independence.

\begin{corollary}\label{cor:nullity-increment}
Under the hypotheses of \Cref{thm:dilation}, hence for every nonsymmetric
dense core, put \(g=g(C)\), and let
\[
 f(q,M)=\#\{r\in\Z/q\Z:2r=M\pmod q\}.
\]
Then
\begin{equation}\label{eq:nullity-increment}
 \dim\K(tC)-\dim\K(C)
 =\frac{(t-1)g-f(tg,tN)+f(g,N)}2.
\end{equation}
If \(g=1\), this increment is \(\lfloor(t-1)/2\rfloor\).
\end{corollary}

\begin{proof}
There are \((t-1)g\) off-grid residue classes modulo \(tg\).
Reflection \(r\mapsto tN-r\) fixes
\(f(tg,tN)-f(g,N)\) of them; the subtracted term counts exactly the fixed
classes on the grid.  Each remaining two-element orbit supports one odd
function.  The residue functions restrict injectively to the available
positions, since every off-grid position is available and every residue
occurs.  This proves \eqref{eq:nullity-increment}.
For \(g=1\), the grid has one fixed residue; reflection modulo \(t\) has
one fixed residue for odd \(t\) and two for even \(t\).
\end{proof}

\begin{example}\label{ex:dilation-ray}
The nonsymmetric dense core \(A=(4,2)\), \(B=(1,0)\) has modulus one and one
available position, namely \(3\).  Its unit is
\(U_3=h_1h_2^2-h_1^2h_3\ne0\).  Thus \(\K(C)=0\), and \Cref{thm:dilation} gives the complete kernel on
every dilation of this core.  Its basis consists of the differences of the
indicator vectors of residue classes \(r\) and \(t-r\), for
\(1\le r<t/2\); all positions divisible by \(t\) have coefficient zero.
The nullity is \(\lfloor(t-1)/2\rfloor\) at every \(t\), although the
cofactors of every dilated matrix remain independent.
\end{example}

\section{Restoration and full equality}\label{sec:full-equality}

Independence of the units, rather than of the cofactors, is what converts the
first-variation kernel into a statement about full skew Schur equality.

\begin{corollary}\label{cor:independent-restorations}
Let \(C\) be a dense zero-free core with \(\K(C)=0\).
For every \(Z\subseteq\Omega_C\), any connected diagram \(D\) with
\(s_D=s_{C[Z]}\) is a translate of \(C[Z]\) or its rotation.
\end{corollary}

\begin{proof}
The lowest ordinary component identifies the residual core polynomial of
\(D\), with its deletion sign, by \eqref{eq:first-components}.
The residual core is connected.  Skew Schur positivity aligns the signs, and
\Cref{thm:nonnegative-rigidity} identifies its literal core with \(C\),
after an optional rotation and translation.  Write its common-beta set as
\(W\).  Equality of the next components gives \(\mathbf1_Z-\mathbf1_W\in\K(C)\).
Independence forces \(Z=W\).  No property of the unknown mate was assumed
beyond connectedness and equality.
\end{proof}

For a nonsymmetric dense core of modulus one with independent units,
\Cref{cor:nullity-increment} shows that its double also has independent
units.  \Cref{cor:independent-restorations} therefore proves rigidity of
every restoration of both cores.  At larger dilation factors the arithmetic
summand may be nonzero; then the next component restricts possible mates but
does not by itself determine their full functions.

\begin{example}\label{ex:unequal-restorations}
Consider the nonsymmetric dense core \(A=(20,10)\), \(B=(5,0)\) and the two
common-beta sets
\[
 Z=\{1,2,6,11,16\},\qquad
 W=\{2,4,9,14,19\}.
\]
This is the fivefold dilation of \Cref{ex:dilation-ray}.
The difference \(\mathbf1_Z-\mathbf1_W\) is the arithmetic relation
exchanging residues one and four modulo five.  Their deletion signs agree,
so \eqref{eq:first-components} gives equal ordinary components of degrees
two and three.  Nevertheless,
\begin{equation}\label{eq:restoration-separator}
 [h_2h_5h_9^2]s_{C[Z]}=-1,\qquad
 [h_2h_5h_9^2]s_{C[W]}=0.
\end{equation}

Here is a direct coefficient calculation.  For \(Z\), the full lists are
\[
 X=(20,16,11,10,6,2,1),\quad
 Y=(16,11,6,5,2,1,0).
\]
Set all positive-index generators except \(h_2,h_5,h_9\) to zero.
The only surviving assignment has column word
\((2,1,5,4,3,7,6)\), in one-based indices, and labels
\((9,0,9,5,0,2,0)\).  Its sign is negative.
For \(W\), the column list is \((19,14,9,5,4,2,0)\), so the first row
of the specialized matrix is zero.  This proves
\eqref{eq:restoration-separator} directly.
\end{example}

\section{Concluding remarks}

The cofactor theorem is unconditional on the dense class, including every
repeated-label case.  The dilation theorem carries a cofactor hypothesis that
the first theorem supplies on the nonsymmetric dense class, so there it is
unconditional too.  On symmetric dense cores that hypothesis does not remain
open but fails, since \Cref{thm:cofactor-relations} gives a nonzero relation
span there; outside the dense class it is an assumption.

The symmetric dense kernel is determined instead by
\Cref{prop:symmetric-unit-kernel}, which
applies to every dilation as well, dilation preserving symmetry and density;
its off-grid part is spanned by the differences \(e_w-e_{tN-w}\) over the
positions not divisible by \(t\).  The containment of \Cref{thm:dilation} can
therefore be strict: for \(A=(4,3)\), \(B=(1,0)\) one has
\(\mathcal A_{\mathrm{off}}(2C)=0\) while
\(U_1^{2C}=U_7^{2C}=h_1(h_5h_6-h_4h_7)\), a relation inherited from neither
summand.

This paper does not classify the \emph{nonsymmetric} dense cores with
vanishing base unit kernel, the hypothesis of
\Cref{cor:independent-restorations}, although \Cref{ex:dilation-ray} exhibits
one; on symmetric cores \Cref{prop:symmetric-unit-kernel} settles it, the
kernel vanishing exactly when \(|\Omega_C|\le1\).  Such a classification would
turn the direct-sum formula into a complete numerical computation.  General
necessity, including rigidity of composition-irreducible connected diagrams
outside proved subclasses, remains open.

\begin{question}\label{q:nondense-cores}
Which connected cores that are not dense have no cofactor relations beyond
those forced by rotation?  \Cref{lem:endpoint-normalization} already applies
without density,
so such an extension would immediately enlarge the domain of
\Cref{thm:dilation}.
\end{question}

\begin{question}\label{q:restoration-charts}
How can the surviving restoration ambiguities be converted into compatible
factor-reversal presentations?  Neither the first-variation kernel nor
\Cref{cor:independent-restorations} supplies a nontrivial composition chart.
Which vectors \(\mathbf1_Z-\mathbf1_W\in\K(C)\) lift to genuine skew
equivalences is the precise question, and \Cref{ex:unequal-restorations}
supplies only its negative side.
\end{question}

\section*{Declarations}

{\footnotesize
\noindent\textbf{Funding.} This research received no external funding.

\smallskip
\noindent\textbf{Competing interests.} The author has no relevant interests to disclose.

\smallskip
\noindent\textbf{Data availability.} No datasets were generated or analyzed.

\smallskip
\noindent\textbf{Computational verification.} The proofs do not depend on numerical computation.  Scripts used for verification are available at \url{https://github.com/johnnyshields/papers}.

\smallskip
\noindent\textbf{Generative AI use.} During the preparation of this work, the author used OpenAI's ChatGPT (GPT-6 Astra, GPT-5.6 Sol) and Anthropic's Claude (Opus 5.5, Opus 5) to assist with mathematical research, manuscript drafting, and the development of auxiliary verification code.  After using these tools, the author reviewed and edited as needed.  The author takes full responsibility for the content of the work.
\par}

\begin{thebibliography}{9}

\bibitem{Barekat2009CoincidencesCorrigendum}
F. Barekat, V. Reiner and S. van Willigenburg,
\emph{Corrigendum to ``Coincidences among skew Schur functions''},
Adv. Math. \textbf{220} (2009), no.~5, 1655--1656,
DOI: \href{https://doi.org/10.1016/j.aim.2008.10.010}{10.1016/\allowbreak j.aim.2008.10.010}.

\bibitem{Billera2006DecomposableCompositions}
L.~J. Billera, H.~Thomas and S.~van Willigenburg,
\emph{Decomposable compositions, symmetric quasisymmetric functions and equality of ribbon Schur functions},
Adv. Math. \textbf{204} (2006), no.~1, 204--240,
DOI: \href{https://doi.org/10.1016/j.aim.2005.05.014}{10.1016/\allowbreak j.aim.2005.05.014}.

\bibitem{Macdonald1995SymmetricFunctions}
I.~G. Macdonald,
\emph{Symmetric Functions and Hall Polynomials}, 2nd ed.,
Oxford Mathematical Monographs, Oxford University Press, 1995,
DOI: \href{https://doi.org/10.1093/oso/9780198534891.001.0001}{10.1093/\allowbreak oso/\allowbreak 9780198534891.001.0001}.

\bibitem{McNamara2009Classification}
P.~R.~W. McNamara and S.~van Willigenburg,
\emph{Towards a combinatorial classification of skew Schur functions},
Trans. Amer. Math. Soc. \textbf{361} (2009), no.~8, 4437--4470,
DOI: \href{https://doi.org/10.1090/S0002-9947-09-04683-2}{10.1090/\allowbreak S0002-9947-09-04683-2}.

\bibitem{Reiner2007Coincidences}
V.~Reiner, K.~M. Shaw and S.~van Willigenburg,
\emph{Coincidences among skew Schur functions},
Adv. Math. \textbf{216} (2007), no.~1, 118--152,
DOI: \href{https://doi.org/10.1016/j.aim.2007.05.006}{10.1016/\allowbreak j.aim.2007.05.006}.

\end{thebibliography}

\end{document}
